\section{Experiments}
\label{sec:exp}

\subsection{Setup}

\paragraph{Dataset.}
We evaluate on exhaustive SAT-based sweeps at three input widths.
For $N=3$, the sweep covers all $254$ non-trivial functions; for $N=4$ it
covers $65{,}534$ functions; and for $N=6$ we use a catalog of $53$
non-series-parallel (NSP) functions, which are the hardest instances because
they admit no series-parallel realization.
Each sweep enumerates, per function, one synthesis attempt per transistor
budget, recording the realized network when the budget is satisfiable; the
smallest satisfiable budget is that function's exact optimum.
Every function contributes \emph{two} synthesis targets: the pull-down
network realizing $f$ and the pull-up network realizing $\neg f(\neg x)$.
Reference counts for the pull-up target are obtained by matching the dual
function $\neg f(\neg x)$ back to its own entry in the sweep.
This yields $508$ targets at $N=3$ and $131{,}068$ at $N=4$.
Because the sweeps are exhaustive, evaluation is against exact optima rather
than a held-out split; we additionally freeze a random subset of $256$
solved $4$-input targets as a fixed set for the Phase~3 comparison.

\paragraph{Implementation.}
\method{} uses an MPNN backbone with $d=128$, $L=3$ layers, and MLP hidden
dimension $256$, followed by a self-attention block over the nodes of each
graph before the policy heads.
Phase~1 uses the path-order supervision of Section~\ref{par:order} with one
implementation per function, trained for $200$ epochs with Adam
(cosine-annealed lr $2\times10^{-3}\!\to\!10^{-4}$) and a $2\times$ weight on
the Head-1 term.
Phase~2 uses $K=16$ trajectories per update, lr $10^{-4}$,
$\varepsilon=0.2$, $\lambda_H=0.05$, up to $200$ updates per target with
early stopping when the exact optimum is reached, and two independent
restarts from the pretrained checkpoint.
Reported transistor counts are the reduced counts $T_{\mathrm{eff}}$ of
Section~\ref{sec:reduction}.
Phase~3 initializes from the Phase~2 checkpoint, generates paired PDN+PUN
trajectories per function, and evaluates complete pairs with ASTRAN
(freePDK45, SA quality~1, weights $(3,4,1,4,2)$).
All experiments run on a single NVIDIA H100 GPU.

\begin{table}[t]
\centering
\caption{Benchmark scale. Each function yields a PDN and a PUN target.}
\label{tab:datasets}
\begin{tabular}{lccc}
\toprule
\textbf{Benchmark} & \textbf{Functions} & \textbf{Targets} & \textbf{Optimum} \\
\midrule
$N=3$ sweep      & $254$    & $508$     & $1$--$8$ \\
$N=4$ sweep      & $65{,}534$ & $131{,}068$ & $4$--$12$ \\
$N=6$ NSP        & $53$     & $53$      & $5$--$7$ \\
\bottomrule
\end{tabular}
\end{table}

\paragraph{Baselines.}
We compare against:
\begin{itemize}[leftmargin=*, topsep=2pt, itemsep=0pt]
  \item \textbf{MuSTNet}~\cite{muSTNet}: SAT-based exact synthesis.
  \item \textbf{MiniTNtk}~\cite{miniTNtk}: heuristic transistor-network toolkit.
  \item \textbf{MCTS-PTNS}~\cite{mctsptns}: Monte Carlo Tree Search synthesis.
  \item \textbf{\method{} Phase~1}: NLL pretraining only.
  \item \textbf{\method{} Phase~2}: transistor-count RL (single network).
\end{itemize}
For physical metrics (Phases~2--3), we generate the SPICE netlist for
each synthesized PDN+PUN pair and evaluate it with ASTRAN using the same
placement parameters to ensure a fair comparison.

\paragraph{Metrics.}
\textit{Logical metrics} (Phases 1--2):
\begin{itemize}[leftmargin=*, topsep=2pt, itemsep=0pt]
  \item \textbf{Success rate}: fraction of functions with a correct implementation.
  \item \textbf{Optimality rate}: fraction matching the known minimum transistor count.
  \item \textbf{Avg.\ transistors}: mean transistor count over correct solutions.
\end{itemize}
\textit{Physical metrics} (Phase 3):
\begin{itemize}[leftmargin=*, topsep=2pt, itemsep=0pt]
  \item \textbf{ASTRAN objective} $\Phi$ (Eq.~\ref{eq:phi}): weighted sum of
        CPP, gate mismatches, WL, routing density, and diffusion cuts.
  \item \textbf{CPP} (cell width): total poly columns used by the placed cell.
  \item \textbf{Gate mismatches}: positions where P-row and N-row gate signals differ.
  \item \textbf{Nr.\ Gaps}: number of diffusion cuts (Euler-path breaks).
  \item \textbf{WL / Rt.\ Density}: wire length and peak routing congestion.
\end{itemize}

\subsection{Main Results}

Table~\ref{tab:main} reports optimality rates for the pretrained policy
(Phase~1) and after transistor-count finetuning (Phase~2).
On the exhaustive $3$-input sweep, \method{} attains the exact optimum on
\textbf{all $508$ targets}.
The decomposition is informative: the pretrained policy alone is already
optimal on $97.6\%$ of targets at its first sampled evaluation, and Phase~2
closes the remaining $12$ targets, each within $12$ updates.
This indicates that when supervision is consistent and ordered along
conduction paths, imitation alone recovers exact-optimal synthesis at this
width, and reinforcement learning is needed only for a small residual.

At $N=4$ the residual is substantially larger and Phase~2 carries more of the
load; optimization over the full sweep is ongoing and we report the rate over
completed targets.
At $N=6$ the two phases separate cleanly: the pretrained policy produces a
\emph{correct} network for $50$ of $53$ NSP functions but at a mean reduced
cost of $19.0$ transistors against optima near $7$, so completion is largely
solved by pretraining while compaction is the task left to Phase~2.

\begin{table}[t]
\centering
\caption{Optimality rate (fraction of targets reaching the exact SAT optimum).
         $N{=}4$ is over completed targets; optimization is ongoing.}
\label{tab:main}
\resizebox{\columnwidth}{!}{%
\begin{tabular}{lccc}
\toprule
\textbf{Method} & \textbf{$N{=}3$ (508)} & \textbf{$N{=}4$} & \textbf{$N{=}6$ NSP} \\
\midrule
MuSTNet~\cite{muSTNet}     & $100$ & $100$ & $100$ \\
MiniTNtk~\cite{miniTNtk}   & \TODO{} & \TODO{} & \TODO{} \\
MCTS-PTNS~\cite{mctsptns}  & \TODO{} & \TODO{} & \TODO{} \\
\method{} Phase~1          & $97.6$ & \TODO{} & $0$ \\
\method{} Phase~2 (ours)   & $\mathbf{100}$ & $88^{\dagger}$ & \TODO{} \\
\bottomrule
\end{tabular}}
\end{table}

\noindent
MuSTNet is exact by construction and therefore attains $100\%$ optimality;
its cost is a per-instance SAT invocation, whereas \method{} amortizes
synthesis into a learned policy and requires no solver at inference.
\TODO{Fill MiniTNtk / MCTS-PTNS rows and add an inference-time column once
those baselines are run on the same sweeps.}

\subsection{Ablation: Supervision Order Dominates Architecture}
\label{sec:ablation_order}

Our central design claim is that \emph{how} a reference network is linearized
into a supervision sequence matters more than model capacity.
Table~\ref{tab:ablation_order} varies the two factors independently at
$N=4$.  Under index/breadth-first order, Head-1 accuracy saturates near
$0.54$--$0.63$ regardless of capacity; under the path order of
Section~\ref{par:order} it exceeds $0.88$ with the \emph{same} encoders.
The ordering is worth roughly $+30$ points, while doubling width and adding
attention is worth $4$--$8$.

\begin{table}[t]
\centering
\caption{Head-1 (source-selection) accuracy at $N{=}4$:
         architecture $\times$ supervision order.}
\label{tab:ablation_order}
\begin{tabular}{lcc}
\toprule
\textbf{Architecture} & \textbf{Index/BFS order} & \textbf{Path order} \\
\midrule
$d{=}64$, MLP $128$               & $0.543$ & $0.886$ \\
$d{=}128$, MLP $256$, attention   & $0.625$ & $\mathbf{0.932}$ \\
\bottomrule
\end{tabular}
\end{table}

The effect is not specific to $N=4$: at $N=6$ the same substitution raises
Head-1 accuracy from $0.505$ to $0.91$.
We attribute this to the structural argument of Section~\ref{par:order}: under
the path order the source is determined by the graph during walk extensions,
whereas under an index order it is a convention the model must memorize.

\paragraph{Effect of label consistency.}
Retaining every implementation of a function introduces conflicting targets
for identical states.  At $N=6$, restricting supervision to one
implementation per function raises Head-1 accuracy from $0.443$ to $0.505$;
in the multi-implementation setting the loss also degrades over training
rather than continuing to descend, consistent with an irreducible
label-conflict floor.

\paragraph{Effect of network reduction.}
Because redundant devices are removable in post-processing
(Lemma~\ref{lem:deletion}), scoring unreduced networks penalizes the policy
for structure a downstream pass deletes.  On the $N=6$ NSP benchmark, mean
generated size falls from $41.4$ (unreduced) to $40.1$ (structural reduction
only) to $19.0$ (full irredundant reduction), and one function reaches its
exact optimum under reduction alone.

\paragraph{Effect of the search budget.}
Completion depends strongly on the rollout step budget: at $N=6$, raising the
budget from $40$ to $80$ placements raises the fraction of functions for
which a correct network is ever found from $13/36$ to $30/36$, while mean size
grows, confirming that under a tight budget the policy fails to complete
rather than failing to compact.

\paragraph{Planned Phase-2 ablations.}
\TODO{Add EMA vs.\ GRPO baseline, fixed vs.\ adaptive $t^*$, and entropy
$\lambda_H \in \{0, 0.05\}$ rows; each requires a matched multi-seed run
(see Section~\ref{sec:planned}).}

\subsection{Analysis}
\label{sec:analysis}

\paragraph{Top-$k$ behaviour of the pretrained policy.}
Although top-$1$ Head-1 accuracy is $0.50$ on $6$-input states, the correct
node lies in the top $3$ candidates $78\%$ of the time and in the top $5$
$89\%$ of the time (mean rank $1.6$ of ${\sim}64$).
Since Phase~2 \emph{samples} from the policy rather than taking its argmax,
top-$k$ mass — not top-$1$ accuracy — is the relevant measure of prior
quality, which explains why a policy with $0.50$ top-$1$ accuracy still
finetunes effectively.

\paragraph{Where source selection fails.}
Accuracy is $0.62$ on small partial networks but falls to $0.36$ on
mid-sized ones, and is markedly lower when the target is an interior node
($0.37$) than a rail ($0.66$), with the model over-predicting rails.
Errors are largely not near-misses and the model is confidently wrong rather
than undecided, indicating a learning rather than a tie-breaking failure and
motivating the node-feature extension in Section~\ref{sec:planned}.

\paragraph{Representational ceiling.}
A Weisfeiler--Leman refinement over the unified graph states assigns every
candidate node a distinct colour within two rounds, so the achievable Head-1
accuracy is $1.0$.
The observed gap is therefore an optimization and inductive-bias limitation,
not a limitation of the state encoding.

\paragraph{Cost of the reward variants.}
All Phase-2 reward variants cost essentially the same per update
($\approx 1.6$\,s under matched conditions), because rollout generation
dominates while the baseline and exploration terms are comparatively free.
Throughput is instead governed by device sharing: per-update time grows from
$1.9$\,s with one process per GPU to $20$\,s at five and $49$\,s at sixteen,
since each generation step issues many small kernels that serialize.
Wall-clock comparisons between variants are therefore only meaningful at
matched concurrency.

\subsection{Phase 3: Physical-Aware Placement Results}
\label{sec:physical_exp}

\paragraph{Setup.}
For each method, we generate the SPICE netlist for the best PDN+PUN pair
found over 20 trials and evaluate it with ASTRAN (freePDK45, same placement
parameters as training).
The physical RL baseline starts from the Phase~2 checkpoint and is finetuned
for \TODO{} epochs.
Unlike Phase~2, where each network is a separate target, the physical
objective is only defined once \emph{both} the pull-down and pull-up networks
are complete, since ASTRAN places the combined cell; a Phase~3 target is
therefore a \emph{function}, and the reward is assigned jointly to the two
trajectories.
\TODO{Phase~3 numbers pending: results below are to be filled once the
physical evaluation runs over the frozen $256$-target $4$-input set and the
full $3$-input set.}

\begin{table}[t]
\centering
\caption{Physical placement comparison on XOR-3 (20 generation trials).
         $\Phi$ is the ASTRAN weighted objective (Eq.~\ref{eq:phi});
         lower is better.}
\label{tab:physical}
\begin{tabular}{lccccc}
\toprule
\textbf{Method} & $\Phi$ & \textbf{CPP} & \textbf{GM} & \textbf{WL} & \textbf{Gaps} \\
\midrule
\method{} Phase~2 (no physical RL) & \TODO{} & \TODO{} & \TODO{} & \TODO{} & \TODO{} \\
\method{} Phase~3 (ours)           & \TODO{} & \TODO{} & \TODO{} & \TODO{} & \TODO{} \\
\bottomrule
\end{tabular}
\end{table}

\paragraph{Physical RL reduces placement cost.}
\TODO{Discuss: Phase~3 reduces $\Phi$ by X\% vs Phase~2 on XOR-3.
Key gains: CPP drops from Y to Z (fewer poly columns), gate mismatches
reduced to near-zero.  Nr.\ Gaps improvement quantifies better dual
Euler path quality.  Concrete result from training log: best\_obj
improved from 159 (Phase~2 baseline on NET.sp) to \TODO{}.}

\paragraph{CPP and Width relationship.}
For a cell with $T_{\mathrm{PDN}}$ and $T_{\mathrm{PUN}}$ transistors,
cell width satisfies:
\[
  \mathrm{CPP} = 2 \cdot \max(T_{\mathrm{PDN}}, T_{\mathrm{PUN}}) + 1
               + (\text{\# unique gap positions}).
\]
For our XOR-3 result (8+8 transistors), the minimum achievable CPP is 17
(perfect dual Euler path, zero gaps); Phase~3 achieves CPP=18 with
Nr.\ Gaps=2, just one column above the theoretical minimum.

\subsection{Generalization to Unseen Functions}

\TODO{Train Phase~1 on a subset of the $N{=}4$ sweep and evaluate on the
held-out remainder without retraining. Note that the main results above are
evaluated against exhaustive optima rather than a held-out split, so this
experiment is what isolates generalization from memorization.}

\subsection{Scaling Analysis}

\TODO{Report inference time per function as a function of $N$ against the
SAT baseline. The relevant claim is amortization: \method{} pays a one-time
training cost and a fixed sampling cost per function, whereas exact methods
pay a per-instance solver cost that grows with $N$.}

\subsection{Planned Experiments}
\label{sec:planned}

The following studies are required to complete the evaluation.

\paragraph{Necessity of pretraining.}
Running Phase~2 from a randomly initialized policy, versus from the Phase~1
checkpoint, isolates the contribution of imitation pretraining and validates
the two-phase design.

\paragraph{Variance across seeds.}
Per-target outcomes on hard instances are stochastic: in matched single runs
we have observed reward variants exchange wins on the \emph{same} instance.
All variant comparisons will therefore be repeated across seeds and reported
with dispersion rather than as single runs; the numbers in
Section~\ref{sec:ablation_order}, which concern pretraining accuracy over
tens of thousands of targets, are not affected by this.

\paragraph{Exploration-mechanism ablation.}
For instances where the policy converges one transistor above the optimum,
we add novelty and record-beating bonuses, down-weight self-imitation while
the record is suboptimal, and re-heat the sampling temperature on stalls.
A leave-one-out study is needed to attribute the gain to individual
mechanisms rather than to the bundle.

\paragraph{Node features.}
The failure analysis in Section~\ref{sec:analysis} localizes errors to
interior nodes of mid-sized networks.  Supplying per-node structural and
functional descriptors --- distance to the rails, degree, and the number of
uncovered on-patterns routed through each node --- targets exactly this
deficit.

\paragraph{Remaining baselines.}
MiniTNtk and MCTS-PTNS must be run on the same sweeps, together with the
naive sum-of-products series-parallel construction as a no-learning floor.
